\subsubsection{Multiple Filters and Output Feature Maps}\label{sec:multiple-filters-and-output-feature-maps}

In the previous section, we learned that for one filter, there is one output feature map. Now suppose the input is still a tiny image of size $64 \times 64 \times 3$, but instead of using \textbf{one filter}, we use \textbf{eight filters}. \emph{What happens?}

\highspace
\begin{flushleft}
    \textcolor{Green3}{\faIcon{question-circle} \textbf{Why do we use multiple filters?}}
\end{flushleft}
Imagine trying to describe a face. If we ask only ``\emph{is there a vertical edge?}'', we obtain only one type of information. But a face contains \emph{eyes}, \emph{nose}, \emph{mouth}, \emph{hair}, \emph{contours}, \emph{textures}, and so on. One detector cannot capture everything.

\highspace
Therefore, CNNs learn \textbf{many different filters simultaneously}. Precisely, \hl{different filters specialize in different visual patterns.}

\highspace
\textcolor{Green3}{\faIcon{book} \textbf{Every filter learns a different pattern.}} Suppose we have four filters:
\begin{itemize}
    \item Filter 1 might respond to \emph{vertical edges}.
    \item Filter 2 might respond to \emph{horizontal edges}.
    \item Filter 3 might respond to \emph{corners}.
    \item Filter 4 might respond to \emph{textures}.
\end{itemize}
Notice that \textbf{all filters receive exactly the same input volume}. The difference lies only in their weights.

\highspace
\textcolor{Green3}{\faIcon{tachometer-alt} \textbf{All filters operate in parallel.}} Every filter independently performs exactly the computation described in the previous section (\autopageref{sec:convolution-over-multi-channel-inputs}, precisely, the flow described on \autopageref{fig:convolution-over-multi-channel-inputs}).

\highspace
For example, suppose the input is a $64 \times 64 \times 3$ image, and every filter produces a $62 \times 62 \times 1$ output. Instead of keeping $N$ maps separate, we \textbf{stack} them. If $N = 4$, the depth of the output volume is therefore \textbf{the number of filters}. In this case, the output volume is $62 \times 62 \times 4$.

\highspace
\textcolor{Green3}{\faIcon{question-circle} \textbf{Why does the depth increase?}} On page \autopageref{sec:depth-increases-while-width-decreases}, we observed an interesting pattern in CNNs: as we \hl{move deeper into the network, the number of channels (depth) generally increases, while the spatial dimensions (width and height) decrease.} We can now \textbf{understand one of the reasons} why this happens. The depth increases because a convolutional layer typically \textbf{applies multiple filters}, and \textbf{each filter produces exactly one output feature map}. Therefore, the more filters we use, the deeper the output volume becomes. This also explains why deeper CNNs are able to represent increasingly rich visual information: by learning many different filters, each layer can specialize in detecting different patterns, from simple edges and textures to progressively more complex structures.

\highspace
\begin{takeawaysbox}
    This section can be summarized by three rules:
    \begin{enumerate}
        \item \textbf{One filter} always produces \textbf{one output} feature map.
        \item \textbf{Different filters} learn \textbf{different visual patterns} from the same input volume.
        \item The output \textbf{depth} of a convolutional layer is exactly the \textbf{number of filters}.
    \end{enumerate}
\end{takeawaysbox}